\documentclass[12pt,a4paper]{article}
\usepackage[margin=25mm, headheight=15pt, headsep=10pt]{geometry}
\usepackage{fontspec}
\usepackage[table]{xcolor} % Note: table option is required for rowcolors
\usepackage{titlesec}
\usepackage{tocloft}
\usepackage{fancyhdr}
\usepackage{booktabs}
\usepackage{tcolorbox}
\tcbuselibrary{skins, breakable}
\usepackage[colorlinks=true, linkcolor=emerald, urlcolor=emerald, bookmarks=true]{hyperref}

\definecolor{emerald}{RGB}{0,102,51}
\definecolor{darkred}{RGB}{139,0,0}
\definecolor{lightgray}{RGB}{245,245,245}

\usepackage[nil]{babel}
\babelprovide[import, main]{bengali}
\babelprovide[import]{arabic}
\babelfont{rm}[Renderer=HarfBuzz,Scale=1.0]{Noto Serif Bengali}
\babelfont{sf}[Renderer=HarfBuzz]{Noto Sans Bengali}
\babelfont{tt}[Renderer=HarfBuzz]{Noto Sans Bengali}
\babelfont[arabic]{rm}[Renderer=HarfBuzz,Scale=1.1]{Noto Naskh Arabic}

% Section styling
\titleformat{\section}{\Large\bfseries\color{emerald}}{\thesection.}{0.5em}{}
\titleformat{\subsection}{\large\bfseries\color{emerald!85!black}}{\thesubsection.}{0.5em}{}
\titleformat{\subsubsection}{\normalsize\bfseries\color{emerald!70!black}}{\thesubsubsection.}{0.5em}{}
\titlespacing*{\section}{0pt}{18pt}{8pt}

% Fancy Header/Footer setup
\pagestyle{fancy}
\fancyhf{} % clear all header and footer fields
\fancyhead[L]{\color{emerald}\small\textbf{\nouppercase{\leftmark}}}
\fancyfoot[L]{\scriptsize\color{black!50}{\lat{AI}}-সংকলিত, আলেম-যাচাইকৃত নয়}
\fancyfoot[C]{\color{emerald!70!black}\thepage}
\renewcommand{\headrulewidth}{0.4pt}
\renewcommand{\headrule}{\hbox to\headwidth{\color{emerald}\leaders\hrule height \headrulewidth\hfill}}
\renewcommand{\footrulewidth}{0pt}

% Custom boxes
% 1. Ayat/Hadith Box
\newtcolorbox{ayatbox}[1][]{
    enhanced, breakable, 
    colback=emerald!5, 
    colframe=emerald, 
    boxrule=0.5pt, 
    arc=3pt, 
    left=10pt, right=10pt, top=10pt, bottom=10pt,
    #1
}

% 2. Quote/Translation Box
\newtcolorbox{quotebox}[1][]{
    enhanced, breakable,
    frame hidden,
    colback=lightgray,
    borderline west={3pt}{0pt}{emerald},
    left=15pt, right=10pt, top=10pt, bottom=10pt,
    sharp corners,
    #1
}

% 3. AI-generated notice box
\newtcolorbox{warnbox}[1][]{
    enhanced, breakable,
    colback=red!4,
    colframe=darkred,
    boxrule=0.8pt,
    arc=3pt,
    left=10pt, right=10pt, top=10pt, bottom=10pt,
    #1
}

% Commands
\newcommand{\arab}[1]{\foreignlanguage{arabic}{#1}}
\newcommand{\kib}[1]{\textcolor{emerald}{\textbf{#1}}}

% Latin (English) fallback: Noto Serif Bengali has NO Latin glyphs
\newfontfamily\latfont[Renderer=HarfBuzz]{Noto Serif}
\newcommand{\lat}[1]{{\latfont #1}}

\setlength{\parskip}{5pt}
\setlength{\parindent}{0pt}

% Fix for missing bullet in Noto Serif Bengali
\renewcommand{\labelitemi}{$\bullet$}



\begin{document}
\begin{center}{\Large\bfseries\color{emerald} ঘটনা ৪: নবীজি (সা.) — চাটাইয়ে ঘুম, দুনিয়ায় মুসাফিরের মতো}\end{center}
\vspace{1em}
\section*{ঘটনা ৪: নবীজি (সা.) — চাটাইয়ে ঘুম, দুনিয়ায় মুসাফিরের মতো}

\begin{warnbox}
 \textbf{গুরুত্বপূর্ণ নোটিশ (\lat{Important} \lat{Notice}):} এই কন্টেন্টটি \lat{AI} (কৃত্রিম বুদ্ধিমত্তা) দিয়ে প্রস্তুত ও সংকলিত। \lat{AI} কঠোর সূত্র-মান নীতি মেনে শুধু নির্ভরযোগ্য উৎস থেকে তথ্য সংগ্রহ করেছে — কেবল সহীহ/হাসান হাদিস ও সহীহ সনদের আসার রাখা হয়েছে, দুর্বল সনদ বাদ দেওয়া হয়েছে। তবে এটি কোনো আলেম বা মুহাদ্দিস কর্তৃক যাচাইকৃত নয়।
\end{warnbox}


\medskip\hrule\medskip

\subsection*{১. সূত্র কার্ড}

\begin{table}[h]\centering\small
\begin{tabular}{ll}
বিষয় & বিবরণ \\
\hline
\textbf{বর্ণনাকারী} & আবদুল্লাহ ইবনে মাসউদ (রাঃ) থেকে \\
\textbf{গ্রন্থ} & সুনানে তিরমিজি, হাদিস ২৩৭৭; সহীহ মুসলিম ২৮৬৩ (অনুরূপ) \\
\textbf{অধ্যায়} & কিতাবুজ-যুহদ, বাব: দুনিয়া সম্পর্কে নবীজির অবস্থান \\
\textbf{সনদের মান} & \textbf{হাসান সহীহ} (তিরমিজির ঘোষণা) \\
\textbf{মূল বিষয়} & দুনিয়ার প্রতি অনাসক্তি (যুহদ) — অহংকার ও দুনিয়ামোহ থেকে মুক্ত থাকার উৎস \\
\hline
\end{tabular}
\end{table}

\medskip\hrule\medskip

\subsection*{২. পটভূমি (কে, কখন, কোথায়, কেন)}

মদীনায় ইসলামি রাষ্ট্র প্রতিষ্ঠার পরও নবীজি (সাল্লাল্লাহু আলাইহি ওয়া সাল্লাম) অত্যন্ত \lat{simple} (সাধারণ) জীবনযাপন করতেন। তাঁর ঘরে প্রায়ই সামান্য \lat{belongings} (জিনিস) থাকত। সাহাবারা তাঁকে ভালোবেসে বলতেন — \lat{“}আপনার জন্য আরামদায়ক বিছানা বানিয়ে দিই।\lat{”}

একদিন নবীজি একটি \textbf{\lat{mattress} (চাটাই)} (খেজুর পাতার তৈরি মোটা বুনন) শুয়েছিলেন। চাটাইয়ের শক্ত বুননের \lat{marks} (দাগ) তাঁর পবিত্র শরীরে বসে গিয়েছিল। সাহাবারা সেটি দেখে মন খারাপ করলেন — এবং সেই ভালোবাসা থেকেই নবীজিকে একটি \lat{comfortable} (আরামদায়ক) বিছানার প্রস্তাব দিলেন।

নবীজির \lat{reply} (উত্তর) টি ছিল এমন — যা সাহাবাদের চোখে পানি আনে এবং উম্মতের জন্য দুনিয়ার প্রতি \lat{attitude} (মনোভাব)-এর চিরন্তন \lat{lesson} (শিক্ষা) হয়ে থাকে।

\medskip\hrule\medskip

\subsection*{৩. পূর্ণ ঘটনার বর্ণনা}

আবদুল্লাহ ইবনে মাসউদ (রাঃ) বলেন:

\begin{quote}
\textbf{\lat{“}নবীজি (সাল্লাল্লাহু আলাইহি ওয়া সাল্লাম) একদিন একটি চাটাইয়ে ঘুমালেন। উঠলে চাটাইয়ের দাগ তাঁর (পবিত্র) পাশে পড়ে গেছে।} \\
\textbf{আমরা বললাম: "হে আল্লাহর রাসূল! আমরা আপনার জন্য যদি বিছানা (আরামদায়ক ও নরম) বানিয়ে দিই?"} \\
\textbf{নবীজি (সা.) বললেন: "দুনিয়ার সাথে আমার কী সম্পর্ক? আমি তো দুনিয়ায় এমন এক আরোহীর মতো, যে একটি গাছের ছায়ায় একটু বিশ্রাম নিল, তারপর রওনা হয়ে গাছটি ছেড়ে চলে গেল।"\lat{”}}
\end{quote}

\medskip\hrule\medskip

\subsection*{৪. শব্দের ব্যাখ্যা}

\begin{table}[h]\centering\small
\begin{tabular}{ll}
শব্দ & অর্থ \\
\hline
\textbf{হাসীর (\arab{حصير})} & চাটাই — খেজুর পাতা বা নলখাগড়ার তৈরি শক্ত মোটা বুনন \\
\textbf{উইতাআ (\arab{وطاء})} & নরম বিছানা/আসন \\
\textbf{রাকিব (\arab{راكب})} & আরোহী — সফরে থাকা মুসাফির \\
\textbf{যুহদ (\arab{زهد})} & দুনিয়ার প্রতি অনাসক্তি; প্রয়োজন অনুযায়ী গ্রহণ, কিন্তু দুনিয়াকে অন্তরে প্রাধান্য না দেওয়া \\
\hline
\end{tabular}
\end{table}

\textbf{গাছের ছায়ার উপমা:} পথিক গরমে গাছতলায় ছায়া নেয় — কিন্তু সে জানে ছায়া স্থায়ী নয়, তাকে পথ চলতেই হবে। ঠিক তেমনি দুনিয়া — সাময়িক; আখেরাতই স্থায়ী গন্তব্য। এই চেতনা থাকলে দুনিয়ার অহংকার, বিলাসিতা ও গর্ব মনে আসে না।

\medskip\hrule\medskip

\subsection*{৫. সংশ্লিষ্ট হাদিস ও আয়াত}

\textbf{হাদিস (বুখারি ২৩৭৭-এর অনুরূপ — মুসলিম ২৮৬৩):} নবীজি বলেছেন: \textbf{\lat{“}দুনিয়া হলো মুমিনের কারাগার ও কাফিরের জান্নাত।\lat{”}}

\textbf{আয়াত (সূরা আল-হাদীদ ৫৭:২০):}
\begin{quote}
\textbf{\lat{“}জেনে রাখো, দুনিয়ার জীবন তো খেল-তামাশা, অলংকার, তোমাদের পারস্পরিক গর্ব ও ধন-সন্তানে আধিক্য ছাড়া আর কিছুই নয়।\lat{”}}
\end{quote}

\textbf{আয়াত (সূরা আল-কাহফ ১৮:৪৫):}
\begin{quote}
\textbf{\lat{“}তাদের কাছে দুনিয়ার জীবনের উপমা বর্ণনা করুন — তা পানির মতো, যা আমি আকাশ থেকে বর্ষণ করি... অতঃপর তা শুষ্ক তুষের মতো হয়ে যায়, যা বাতাস উড়িয়ে নিয়ে যায়।\lat{”}}
\end{quote}

\textbf{আয়াত (সূরা আয-যুখরুফ ৪৩:৩৫):}
\begin{quote}
\textbf{\lat{“}আর এ সবই তো দুনিয়ার জীবনের সাময়িক ভোগের উপকরণ। আর আখেরাত আপনার রবের কাছে মুত্তাকীদের জন্য।\lat{”}}
\end{quote}

\textbf{হাদিস (তিরমিজি ২৩৭৮):} নবীজি বলেছেন: \textbf{\lat{“}যে দুনিয়াকে ভালোবাসবে, তার আখেরাত ক্ষতিগ্রস্ত হবে; আর যে আখেরাতকে ভালোবাসবে, তার দুনিয়া ত্যাগের প্রস্তুতি হবে।\lat{”}}

\medskip\hrule\medskip

\subsection*{৬. রেফারেন্স}

\begin{itemize}
\item সুনানে তিরমিজি, হাদিস ২৩৭৭ (হাসান সহীহ)
\item সহীহ মুসলিম, হাদিস ২৮৬৩ (দুনিয়া মুমিনের কারাগার)
\item সহীহ বুখারি, হাদিস ৬৪১১ (নবীজির ঘরের সরলতা)
\item সূরা আল-হাদীদ, আয়াত ২০; সূরা আল-কাহফ, আয়াত ৪৫; সূরা আয-যুখরুফ, আয়াত ৩৫
\end{itemize}

\end{document}
